% This is text associated with the open textbook "Open Electromagnetics"
% See immediately below for author, date
% This work is licensed under the Creative Commons Attribution-ShareAlike 4.0 International License. (CC BY-SA 4.0) 
% To view a copy of the license, visit http://creativecommons.org/licenses/by-sa/4.0/

%ORB TITLE Phase and Group Velocity
%ORB AUTHOR 0        # S.W. Ellingson, Virginia Tech (ellingson@vt.edu)
%ORB DATE 20191121   # yyyymmdd
%ORB ABSTRACT 

\label{m0176_Phase_and_Group_Velocity} % <-- Must match name of this file and module directory name.  Don't mess with this!
\index{phase velocity|(}
\index{group velocity|(}

Phase velocity is the speed at which a point of constant phase travels as the wave propagates.\footnote{Formally, ``velocity'' is a vector which indicates both the direction and rate of motion.  It is common practice to use the terms ``phase velocity'' and ``group velocity'' even though we are actually referring merely to rate of motion.  The direction is, of course, in the direction of propagation.}
For a sinusoidally-varying wave, this speed is easy to quantify.
To see this, consider the wave:
%
\begin{equation}
%w(z,t) = 
A\cos\left(\omega t -\beta z + \psi \right)
\label{m0176_ewzt}
\end{equation}
%
where $\omega=2\pi f$ is angular frequency, $z$ is position, and $\beta$ is the phase propagation constant.
At any given time, the distance between points of constant phase is one wavelength $\lambda$.
Therefore, the phase velocity $v_p$ is
%
\begin{equation}
v_p = \lambda f 
\label{m0176_evp}
\end{equation}
%
Since $\beta=2\pi/\lambda$, this may also be written as follows:
%
\begin{equation}
\boxed{
v_p = \frac{\omega}{\beta} 
}
\label{m0176_evp2}
\end{equation}
%
Noting that $\beta=\omega\sqrt{\mu\epsilon}$ for simple matter, we may also express $v_p$ in terms of the constitutive parameters $\mu$ and $\epsilon$ as follows:
%
\begin{equation}
v_p = \frac{1}{\sqrt{\mu\epsilon}}
\label{m0176_evpm}
\end{equation}
%
Since $v_p$ in this case depends \emph{only} on the constitutive properties $\mu$ and $\epsilon$, it is reasonable to view phase velocity also as a property of matter.

Central to the concept of phase velocity is uniformity over space and time.
Equations~\ref{m0176_evp}--\ref{m0176_evpm} presume a wave having the form of Equation~\ref{m0176_ewzt}, which exhibits precisely the same behavior over all possible time $t$ from $-\infty$ to $+\infty$ and over all possible $z$ from $-\infty$ to $+\infty$.
This uniformity over all space and time precludes the use of such a wave to send information.
To send information, the source of the wave needs to vary at least one parameter as a function of time; for example 
$A$ (resulting in amplitude modulation), $\omega$ (resulting in frequency modulation), or $\psi$ (resulting in phase modulation).
In other words, information can be transmitted only by making the wave non-uniform in some respect.
Furthermore, some materials and structures can cause changes in $\psi$ or other combinations of parameters which vary with position or time.  
Examples include dispersion
\index{dispersion}
and propagation within waveguides.
Regardless of the cause, varying the parameters $\omega$ or $\psi$ as a function of time means that the instantaneous distance between points of constant phase may be very different from $\lambda$.
Thus, the \emph{instantaneous} frequency of variation as a function of time and position may be very different from $f$.
In this case Equations~\ref{m0176_evp}--\ref{m0176_evpm} may not necessarily provide a meaningful value for the speed of propagation.
   
Some other concept is required to describe the speed of propagation of such waves.
That concept is \emph{group velocity}, $v_g$, defined as follows:
%%%%%%%%%%%%%%%%%%%%%%%%%%%%%%%%%%%%%%%%%%%%%%%
%ORB HIGHLIGHT
\begin{mdframed}[backgroundcolor=yellow!20,nobreak=true] 
\emph{Group velocity}, $v_g$, is the ratio of the apparent change in frequency $\omega$ to the associated change in the phase propagation constant $\beta$; i.e., $\Delta\omega/\Delta\beta$.
\end{mdframed}
%%%%%%%%%%%%%%%%%%%%%%%%%%%%%%%%%%%%%%%%%%%%%%%
Letting $\Delta\beta$ become vanishingly small, we obtain
%
\begin{equation}
\boxed{
v_g \triangleq \frac{\partial\omega}{\partial\beta}
}
\label{m0176_evg}
\end{equation}
%
Note the similarity to the definition of phase velocity in Equation~\ref{m0176_evp2}.
Group velocity can be interpreted as the speed at which a \emph{disturbance} in the wave propagates.
Information may be conveyed as meaningful disturbances relative to a steady-state condition, so group velocity is also the speed of information in a wave.

Note Equation~\ref{m0176_evg} yields the expected result for waves in the form of Equation~\ref{m0176_ewzt}:
%
\begin{flalign}
v_g &= \left(\frac{\partial\beta}{\partial\omega}\right)^{-1} = \left(\frac{\partial}{\partial\omega}\omega\sqrt{\mu\epsilon}\right)^{-1} \nonumber \\
    &= \frac{1}{\sqrt{\mu\epsilon}} = v_p
\end{flalign}
%
In other words, the group velocity of a wave in the form of Equation~\ref{m0176_ewzt} is equal to its phase velocity. 

To observe a difference between $v_p$ and $v_g$,  $\beta$ must somehow vary as a function of something other than just $\omega$ and the constitutive parameters.  Again, modulation (introduced by the source of the wave) and dispersion (frequency-dependent constitutive parameters) are examples in which $v_g$ is not necessarily equal to $v_p$.   
Here's an example involving dispersion:

%%%%%%%%%%%%%%%%%%%%%%%%%%%%%%%%%%%%%%%%%%%%%%%
%ORB EXAMPLE
~\\ \begin{example} 
\begin{mdframed}[backgroundcolor=brown!20,linecolor=brown!20]\raggedright
Phase and group velocity for a material exhibiting square-law dispersion.
\\~\\
A broad class of non-magnetic dispersive media exhibit relative permittivity $\epsilon_r$ that varies as the square of frequency over a narrow range of frequencies centered at $\omega_0$.  For these media we presume
%
\begin{equation}
\epsilon_r = K\left(\frac{\omega}{\omega_0}\right)^2
\end{equation}
% 
where $K$ is a real-valued positive constant.
What is the phase and group velocity for a sinusoidally-varying wave in this material?

{\bf Solution.}
First, note 
%
\begin{flalign}
\beta &= \omega\sqrt{\mu_0\epsilon} = \omega\sqrt{\mu_0\epsilon_0}\sqrt{\epsilon_r} \nonumber \\
      &= \frac{\sqrt{K}\cdot\omega^2}{\omega_0}\sqrt{\mu_0\epsilon_0} \label{m0176_eex1beta}
\end{flalign}
%
The phase velocity is:
%
\begin{equation}
v_p = \frac{\omega}{\beta} = \frac{\omega_0}{\sqrt{K}\cdot\omega \sqrt{\mu_0\epsilon_0}}
\end{equation}
%
Whereas the group velocity is:
%
\begin{flalign}
v_g &= \frac{\partial\omega}{\partial\beta} = \left( \frac{\partial\beta}{\partial\omega} \right)^{-1} \\
    &= \left( \frac{\partial}{\partial\omega} \frac{\sqrt{K}\cdot\omega^2}{\omega_0}\sqrt{\mu_0\epsilon_0}  \right)^{-1} \\
    &= \left( 2\frac{\sqrt{K}\cdot\omega}{\omega_0}\sqrt{\mu_0\epsilon_0}  \right)^{-1} 
\end{flalign}
%
Now simplifying using Equation~\ref{m0176_eex1beta}:
%
\begin{flalign}
v_g &= \left( 2\frac{\beta}{\omega} \right)^{-1} \\
    &= \frac{1}{2}\frac{\omega}{\beta} \\
    &= \frac{1}{2}v_p
\end{flalign}
%
Thus, we see that in this case the group velocity is always half the phase velocity.
\end{mdframed}
% note figures will need to go between "\end{mdframed}" and "\end{example}"
\end{example}
%%%%%%%%%%%%%%%%%%%%%%%%%%%%%%%%%%%%%%%%%%%%%%%

Another commonly-encountered example for which $v_g$ is not necessarily equal to $v_p$ is the propagation of guided waves; e.g., waves within a waveguide.
In fact, such waves may exhibit phase velocity greater than the speed of light in a vacuum, $c$. 
However, the group velocity remains less than $c$, which means the speed at which information may propagate in a waveguide is less than $c$.
No physical laws are violated, since the universal ``speed limit'' $c$ applies to information, and not simply points of constant phase.
(See ``Additional Reading'' at the end of this section for more on this concept.)

%%%%%%%%%%%%%%%%%%%%%

{\bf Additional Reading:}
\begin{itemize}
\item ``\href{https://en.wikipedia.org/wiki/Group_velocity}{Group velocity}'' on Wikipedia.
\item ``\href{https://en.wikipedia.org/wiki/Phase_velocity}{Phase velocity}'' on Wikipedia.
\item ``\href{https://en.wikipedia.org/wiki/Speed_of_light}{Speed of light}'' on Wikipedia.
\end{itemize}

\index{phase velocity|)}
\index{group velocity|)}


